% PAGES: 111-112

% This file continues the proof of Lemma 3.3, begun on PDF page 110
% (sec03_c2_1.tex, same agent). \begin{proof} was opened there.

\begin{proof}
If the derivative security is replicable, then there exists a portfolio
made of the $m$ securities which has the same values at time $\tau$ as the
derivative security, regardless of the state the market is at time $\tau$.
Let $\Theta$ be the $m \times 1$ positions vector of the replicating
portfolio. Then,
\begin{equation}
s_\tau \;=\; \Theta^t M_\tau,
\end{equation}
where $M_\tau$ is the payoff matrix of the $m$ securities at time $\tau$;
cf. (3.9) from Lemma 3.2.

Since the market model is arbitrage--free, it follows from Theorem 3.2 that
the value $s_{t_0}$ at time $t_0$ of the replicable derivative security
must be the same as the value $\Theta^t S_{t_0}$ of the replicating
portfolio at time $t_0$, i.e.,
\begin{equation}
s_{t_0} \;=\; \Theta^t S_{t_0},
\end{equation}
where $S_{t_0}$ is the price vector of the $m$ securities at time $t_0$.

Recall from Theorem 3.1 that
\begin{equation}
S_{t_0} \;=\; M_\tau Q,
\end{equation}
where $Q = (Q_k)_{k=1:n}$ is the column vector of the state prices $Q_k >
0$, $k = 1:n$.

From (3.25), (3.26), and (3.24), we conclude that
\begin{equation}
s_{t_0} \;=\; \Theta^t S_{t_0} \;=\; \Theta^t M_\tau Q \;=\; s_\tau Q \;=\;
\sum_{k=1}^n Q_k s_\tau^{\omega^k},
\end{equation}
which is what we wanted to show; cf. (3.23). Note that, for the last
equality from (3.27), we used the row vector -- column vector
multiplication formula (1.5) and the fact that $s_\tau = (s_\tau^{\omega^1}
\ldots s_\tau^{\omega^n})$ is an $1 \times n$ row vector.
\end{proof}

\section{Complete markets}

\begin{definition}
A one period market model with $m$ securities and $n$ market states at time
$\tau > t_0$ is complete if and only if any cash flow at time $\tau$ can be
replicated using a portfolio made of the $m$ securities.

In other words, the one period market model is complete if and only if,
for any $1 \times n$ payoff vector $C_\tau$, there exists an $m \times 1$
positions vector $\Theta$ such that
\begin{equation}
C_\tau \;=\; \Theta^t M_\tau.
\end{equation}
\end{definition}

\begin{theorem}
A one period market model with $m$ securities and $n$ market states at time
$\tau > t_0$ is complete if and only if the market contains $n$
non--redundant securities, i.e., if and only if the row rank of the payoff
matrix $M_\tau$ is equal to $n$, the number of possible states at time
$\tau$.
\end{theorem}

Note that a consequence of Theorem 3.3 is that a complete market must
contain at least as many securities as possible states in the future, i.e.,
that $m \geq n$, which is an intuitive requirement.

\begin{proof}
Let $M_\tau = \mathrm{row}(S_{j\tau})_{j=1:m}$ be the row form of the
payoff matrix $M_\tau$, where $S_{j\tau}$ is the price vector of asset $j$
at time $\tau$, $j = 1:m$; cf. (3.3). From (1.9), it follows that
\begin{equation}
\Theta^t M_\tau \;=\; \sum_{j=1}^m \Theta_j S_{j\tau}.
\end{equation}

Recall from (3.28) that a market is complete if and only if for any $1
\times n$ payoff vector $C_\tau$ there exists an $m \times 1$ positions
vector $\Theta$ such that $C_\tau = \Theta^t M_\tau$. Then, from (3.29), we
find that the market is complete if and only if
\[
\left\langle S_{1\tau}, S_{2\tau}, \ldots, S_{m\tau} \right\rangle
\;=\;
\left\{ \sum_{j=1}^m \Theta_j S_{j\tau} \ \text{with}\ \Theta_j \in \R, j =
1:m \right\}
\;=\;
\R^{1 \times n},
\]
i.e., if $\R^{1 \times n}$ is the space generated by the vectors
$S_{j\tau}$, $j = 1:m$. This happens if and only if there are exactly $n$
linearly independent vectors among the price vectors $S_{1\tau}, S_{2\tau},
\ldots, S_{m\tau}$ of the $m$ securities, and therefore if and only if
there exist exactly $n$ non--redundant securities.
\end{proof}

\begin{lemma}
A one period market model with $n$ securities and $n$ states at time $\tau
> t_0$ is complete if and only if the payoff matrix $M_\tau$ is
nonsingular.
\end{lemma}

\begin{proof}
From Theorem 3.3, it follows a one period market model with $n$ securities
and $n$ states is complete if and only if the matrix $M_\tau$ has row rank
$n$, i.e., it has $n$ linearly independent rows. Since $M_\tau$ is an $n
\times n$ matrix, it follows that $\rank(M_\tau) = n$, and, from Theorem
1.1, we conclude that the matrix $M_\tau$ is nonsingular.
\end{proof}

\medskip
\noindent\rule{\linewidth}{0.4pt}
\medskip

\noindent\textbf{\textit{One Period Binomial Model Example (Continued):}}\\
The one period binomial model is complete.

To see this, recall from (3.5) that the payoff matrix at time $\tau$
corresponding to one period binomial model is
\[
M_\tau \;=\; \begin{pmatrix} e^{r\delta t} & e^{r\delta t} \\ uS_0 & dS_0
\end{pmatrix}.
\]
From (10.6), we obtain that
\[
\det(M_\tau) \;=\; e^{r\delta t}dS_0 - e^{r\delta t}uS_0 \;=\;
e^{r\delta t}S_0(d - u),
\]
and therefore $\det(M_\tau) \neq 0$, since we assumed that $d < u$.

Then, from Theorem 1.2, it follows that the matrix $M_\tau$ is
nonsingular, and, from Lemma 3.4, we conclude that the one period binomial
model is complete. \hfill$\square$

\medskip
\noindent\rule{\linewidth}{0.4pt}
\medskip
